\begin{pastexam}[label={px:ntu-107-ie-ds-12}]{臺大工工所 107 學年度 第 12 題}
（是非；命題正確填 (A)，錯誤填 (B)；答對得 4 分，答錯倒扣 4 分，至多扣到該大題零分）

In a red-black tree, there can be at most 2 red nodes in any path.
\end{pastexam}

\begin{pastsoln}
答案是 (B)，錯誤。

紅黑樹限制的是「紅色節點不能相鄰」（紅色節點的子節點都是黑色），不是紅色節點的總數（資結書
第~10.3~節）。 一條路徑上紅黑可以交錯出現，紅色的個數可以到黑高度那麼多。

反例：取一棵深度 $0$ 到 $5$ 的滿二元樹（$63$ 個節點，鍵用中序 $1$ 到 $63$ 填，是合法的 BST），
深度 $0, 2, 4$ 的節點塗黑，深度 $1, 3, 5$ 的節點塗紅。 檢查紅黑樹的性質：
\begin{itemize}
\item 根（深度 $0$）是黑的；外部節點 \NIL 都是黑的。
\item 紅色節點在奇數深度，子節點在偶數深度（黑色）或是 \NIL（黑色），沒有紅紅相鄰。
\item 從根到任何一個 \NIL，經過的內部黑色節點都是深度 $0, 2, 4$ 那三個，黑高度相同。
\end{itemize}
所以它是紅黑樹，而從根到任何一片葉子的路徑上有 $3$ 個紅色節點（深度 $1, 3, 5$）。 把層數加大，
紅色節點可以有 $\Theta(\log n)$ 個。

能說的是：任何一條從某節點往下到 \NIL 的路徑上，紅色節點的個數不超過黑色節點的個數，這也是紅黑
樹高度 $\le 2\lg(n+1)$ 的來源（資結書定理~10.4）。
\end{pastsoln}
